\documentclass[10pt,a4paper]{amsart}
\usepackage[margin=19mm]{geometry}
\usepackage{amsmath,amssymb,mathtools}
\usepackage[scr=rsfs,cal=euler]{mathalfa}
\usepackage{xfrac}
\usepackage[unicode,hidelinks,bookmarks=false]{hyperref}
\newcommand{\ie}{i.e.\ }
\newcommand{\cf}{cf.\ }
\newcommand{\resp}{respectively\ }
\newcommand{\Subsection}{\subsection}
\providecommand{\nobreakdash}{\nobreak\hbox{-}}
\setlength{\emergencystretch}{2em}
\multlinegap0pt
\usepackage{tikz}
\usetikzlibrary{cd}
\usepackage{mathtools}
\usepackage{xcolor}
\usepackage{enumerate}
\usepackage{tikz}
\newtheorem{lemma}{Lemma}
\newtheorem{proposition}{Proposition}
\usepackage{cleveref}
\crefname{lemma}{Lemma}{Lemmas}
\crefname{proposition}{Proposition}{Propositions}
\newcommand{\C}{\mathbb{C}}
\newcommand{\p}{\mathbb{P}}
\newcommand{\pain}[1]{\operatorname{P}_{\mathrm{#1}}}
\title{On monodromy of monodromy surfaces}
\author{Pieter Roffelsen, Alexander Stokes}
\date{Source: arXiv:2608.25594v2 (CC BY 4.0)}
\begin{document}
\maketitle

\noindent\emph{Source and licence.} On monodromy of monodromy surfaces. Version arXiv:2608.25594v2 (CC BY 4.0), distributed by arXiv under the Creative Commons Attribution 4.0 International licence, http://creativecommons.org/licenses/by/4.0/, as recorded in the arXiv licence metadata for this exact version and shown on the abstract page https://arxiv.org/abs/2608.25594v2. Full licence notice: \texttt{licenses/arxiv-2608.25594-CC-BY-4.0.txt}.\par\smallskip

\noindent\textit{Editorial scope.} Counted unit: the article's proposition prop:algebraicmonodromyPVI (source0.tex lines 1462--1505). The article's complete original proof of it is delivered with it (source0.tex lines 1507--1740, one \texttt{proof} environment of 234 lines). No mathematics is written, repaired or re-derived: the statement and the proof are the article's own lines, in the article's own notation, and no retained line is substituted or re-flowed.  Retained uncounted as the source-local context the proof needs: lines 1027--1033; lines 1035--1062; lines 1077--1113; lines 1259--1285.  Nothing was omitted from any retained span, so the package carries no omission record.  No retained line contains a citation, so the wrapper carries no bibliography and no bibliography entry.  The retained lines contain the article's own diagram code, which the wrapper typesets with the standard tikz packages; no image file is delivered and the package declares no asset dependency.  Editorial formatting only, in addition: an AMS \texttt{amsart} preamble that restates the article's own declarations and macros used by the retained lines, namely the environments \texttt{lemma}, \texttt{proposition} on separate counters, together with the standard packages those lines need.  The retained environments print in this document as Lemma 1, Proposition 1 in order of appearance, whereas the article prints the same items with the numbering of its own full document.  The complete native source of the article as distributed in the arXiv e-print for this exact version is bundled unchanged as \texttt{sources/208586/source0.tex}.
\begin{equation}\label{eq:linesPVI}
\begin{aligned}
    L_k : \quad x_1 &= b_k+\frac{1}{b_k},\quad &&b_k x_2+x_3=\frac{w_3b_k-w_2}{b_k-b_k^{-1}}, \quad &&k=1,\dots,8, \\
    L_k :\quad  x_2 &= b_k+\frac{1}{b_k},\quad &&b_k x_3+x_1=\frac{w_1b_k-w_3}{b_k-b_k^{-1}}, \quad &&k=9,\dots,16, \\
    L_k : \quad x_3 &= b_k+\frac{1}{b_k},\quad &&b_k x_1+x_2=\frac{w_2b_k-w_1}{b_k-b_k^{-1}}, \quad &&k=17,\dots,24, 
\end{aligned}
\end{equation}

\begin{equation}\label{eq:defbpvi}
\begin{aligned}
    b_1 &= \frac{u_3}{u_2}, &
    b_2 &= \frac{u_2}{u_3}, &
    b_3 &= \frac{u_3u_4}{u_1}, &
    b_4 &= \frac{u_1}{u_3u_4}, &
    b_5 &= u_1, &
    b_6 &= \frac{1}{u_1}, &
    b_7 &= \frac{u_2}{u_4}, &
    b_8 &= \frac{u_4}{u_2},\\
    b_9 &= \frac{u_1}{u_3}, &
    b_{10} &= \frac{u_3}{u_1}, &
    b_{11} &= \frac{u_3u_4}{u_2}, &
    b_{12} &= \frac{u_2}{u_3u_4}, &
    b_{13} &= u_2, &
    b_{14} &= \frac{1}{u_2}, &
    b_{15} &= \frac{u_1}{u_4}, &
    b_{16} &= \frac{u_4}{u_1},\\
    b_{17} &= \frac{u_2}{u_1}, &
    b_{18} &= \frac{u_1}{u_2}, &
    b_{19} &= u_4, &
    b_{20} &= \frac{1}{u_4}, &
    b_{21} &= u_3, &
    b_{22} &= \frac{1}{u_3}, &
    b_{23} &= \frac{u_1u_2}{u_3u_4}, &
    b_{24} &= \frac{u_3u_4}{u_1u_2}.
\end{aligned}
\end{equation}

\begin{lemma}     \label{lem:PvimorphismUtoW}
The definitions of the $u_i,w_i$ in terms of $\theta_j$ yields the field homomorphism 
    \begin{equation} \label{pvifieldextensionUW}
    \C(\mathscr{W}_{\rm{VI}}) \longrightarrow \C(\mathscr{U}_{\rm{VI}}) ,
    \end{equation}
   specified by 
\begin{align*}
	w_1&=u_1+\frac{1}{u_1}+\frac{u_2}{u_3}+\frac{u_3}{u_2}+\frac{u_2}{u_4}+\frac{u_4}{u_2}+\frac{u_1}{u_3 u_4}+\frac{u_3 u_4}{u_1},\\
	w_2&=u_2+\frac{1}{u_2}+\frac{u_1}{u_3}+\frac{u_3}{u_1}+\frac{u_1}{u_4}+\frac{u_4}{u_1}+\frac{u_2}{u_3 u_4}+\frac{u_3 u_4}{u_2},\\
    w_3&=u_3+\frac{1}{u_3}+\hspace{0.5mm}u_4\hspace{0.5mm}+\frac{1}{u_4}+\frac{u_1}{u_2}+\frac{u_2}{u_1}+\frac{u_1 u_2}{u_3 u_4}+\frac{u_3 u_4}{u_1 u_2},\\
	w_4&=4
    +u_3 u_4
    +\frac{1}{u_3 u_4}
    +\frac{u_3}{u_4}
    +\frac{u_4}{u_3}  
    +\frac{u_1}{u_2 u_3}
    +\frac{u_2 u_3}{u_1}
    +\frac{u_1}{u_2 u_4}
    +\frac{u_2 u_4}{u_1}
    +\frac{u_3}{u_1 u_2}
    +\frac{u_1 u_2}{u_3}
    +\frac{u_4}{u_1 u_2}
    +\frac{u_1 u_2}{u_4}\\
    &\quad  +\frac{u_2}{u_1 u_3}
    +\frac{u_1 u_3}{u_2}
    +\frac{u_2}{u_1 u_4}
    +\frac{u_1 u_4}{u_2}
    +\frac{u_1^2}{u_3 u_4}
    +\frac{u_3 u_4}{u_1^2}
    +\frac{u_2^2}{u_3 u_4}
    +\frac{u_3 u_4}{u_2^2}
    +\frac{u_3^2 u_4}{u_1 u_2}
    +\frac{u_1 u_2}{u_3^2 u_4}
    +\frac{u_3 u_4^2}{u_1 u_2}
    +\frac{u_1 u_2}{u_3 u_4^2}.
\end{align*}
 \end{lemma}

\begin{proposition} \label{prop:combinatorialmonodromyPVI}
The combinatorial monodromy of the family $\mathcal{M}\to\mathscr{W}_{\rm{VI}}$ is isomorphic to $W(D_4)$, with an explicit isomorphism 
    \begin{equation*}
     W(D_4)\xrightarrow{\sim}   \operatorname{Fix}_{\operatorname{Aut}(\mathcal{G}_{\rm{VI}})}(\mathcal{G}_{\rm{VI}}^\infty), \quad r\mapsto \tilde{r},
    \end{equation*}
defined by sending the generators $r_k$, $1\leq k\leq 4$, to
\begin{align*}
	\tilde{r}_1&=(3\,\, 7)\,(4\, \, 8)\,(11\,\, 15)\,(12\,\, 16)\,(19\,\, 23)\,(20\,\, 24),\\
	\tilde{r}_2&=(5\,\, 8)\,(6\, \, 7)\,(13\,\, 16)\,(14\,\, 15)\,(21\,\, 24)\,(22\,\, 23),\\
	\tilde{r}_3&=(1\,\, 5)\,(2\, \, 6)\,(9\,\, 14)\,(10\,\, 13)\,(19\,\, 24)\,(20\,\, 23),\\
	\tilde{r}_4&=(3\,\, 8)\,(4\, \, 7)\,(9\,\, 13)\,(10\,\, 14)\,(17\,\, 22)\,(18\,\, 21).
\end{align*}
% \begin{align*}
% 	\tilde{r}_1&=(1\,\, 4)\,(2\, \, 3)\,(9\,\, 12)\,(10\,\, 11)\,(17\,\, 20)\,(18\,\, 19),\\
% 	\tilde{r}_2&=(1\,\, 6)\,(2\, \, 5)\,(9\,\, 14)\,(10\,\, 13)\,(17\,\, 22)\,(18\,\, 21),\\
% 	\tilde{r}_3&=(5\,\, 7)\,(6\, \, 8)\,(9\,\, 11)\,(10\,\, 12)\,(21\,\, 24)\,(22\,\, 23),\\
% 	\tilde{r}_4&=(1\,\, 3)\,(2\, \, 4)\,(13\,\, 15)\,(14\,\, 16)\,(21\,\, 23)\,(22\,\, 24).
% \end{align*}
% {\color{blue} In current numbering/notation as :
% \begin{align*}
% 	\tilde{r}_1&=(3\,\, 7)\,(4\, \, 8)\,(11\,\, 15)\,(12\,\, 16)\,(19\,\, 23)\,(20\,\, 24),\\
% 	\tilde{r}_2&=(5\,\, 8)\,(6\, \, 7)\,(13\,\, 16)\,(14\,\, 15)\,(21\,\, 24)\,(22\,\, 23),\\
% 	\tilde{r}_3&=(1\,\, 5)\,(2\, \, 6)\,(9\,\, 14)\,(10\,\, 13)\,(19\,\, 24)\,(20\,\, 23),\\
% 	\tilde{r}_4&=(3\,\, 8)\,(4\, \, 7)\,(9\,\, 13)\,(10\,\, 14)\,(17\,\, 22)\,(18\,\, 21).
% \end{align*}
% }\todo{AS check if this is the absolutely ideal numbering}
\end{proposition}

\begin{proposition}\label{prop:algebraicmonodromyPVI} 
    % We have the following commutative diagram of homomorphism of fields 
    % \begin{equation}
    %     \begin{tikzcd}
    %         \C(\Gamma^{\rm{VI}})  \arrow[rr] & &\C(\mathscr{U}) \\
    %         & \C (\mathscr{W}^{\rm{VI}}) \arrow[ul, hook, "\rho^*"] \arrow[ur, hook] & 
    %     \end{tikzcd}
    % \end{equation}
    The field extension $\C(\mathscr{W}_{\rm{VI}})\subset \C(\mathscr{U}_{\rm{VI}})$, defined in \Cref{lem:PvimorphismUtoW}, is a Galois extension and the algebraic monodromy of $\mathcal{M}\to\mathscr{W}_{\rm{VI}}$ is 
% of field extensions between $\rho^* \C(\mathscr{W})\subset \C(\Gamma^{\rm{VI}})$ and $\C(\mathscr{W})\subset \C(\mathscr{U})$, so in particular 
    \begin{equation*}
\operatorname{Gal}\left( K_{\Gamma}/\C (\mathscr{W}_{\rm{VI}})\right)\cong \operatorname{Gal}\left(\C(\mathscr{U}_{\rm{VI}})/\C(\mathscr{W}_{\rm{VI}})\right)\cong W(D_4).
    \end{equation*}
    Here, the first isomorphism comes from a $\C(\mathscr{W}_{\rm{VI}})$-linear isomorphism of fields 
    $K_{\Gamma} \to \C(\mathscr{U}_{\rm{VI}})$. 
    % under the natural identification of $\rho^*\C(\mathscr{W}_{\rm{VI}})$ and $\C(\mathscr{W}_{\rm{VI}})$
    % . 
    The second isomorphism comes from the action of $W(D_4)$ on $\C(\mathscr{U}_{\rm{VI}})$ given by 
    \begin{equation} \label{pvi:actionofD4onU}
\begin{aligned}
    r_1 &: &&u_1\to u_1, \quad 
    &&u_2\to u_2, \quad 
    &&u_3\to u_3, \quad 
    &&u_4 \to \frac{u_1 u_2}{u_3 u_4}, \\
    r_2 &: &&u_1\to \frac{u_4}{u_2}, \quad 
    &&u_2\to \frac{u_4}{u_1}, \quad 
    &&u_3\to \frac{u_3 u_4}{u_1 u_2}, \quad 
    &&u_4\to u_4. \\
    r_3 &: &&u_1\to \frac{u_3}{u_2}, \quad 
    &&u_2\to \frac{u_3}{u_1}, \quad 
    &&u_3 \to u_3, \quad 
    &&u_4\to \frac{u_3 u_4}{u_1 u_2}, \\
    r_4 &: &&u_1 \to u_1, \quad 
    &&u_2\to \frac{u_1}{u_3}, \quad 
    &&u_3\to \frac{u_1}{u_2}, \quad
    &&u_4 \to u_4,
\end{aligned}
\end{equation}
which keeps $(w_1,w_2,w_3,w_4)$ fixed and induces the same permutations of lines as in \Cref{prop:combinatorialmonodromyPVI}.
    % and the second isomorphism is given explicitly in terms of an action of $W(D_4)$ on $\C(\mathscr{U}_{\rm{VI}})$.
    % where the field extension $\C(\mathscr{W}_{\rm{VI}})\subset \C(\mathscr{U}_{\rm{VI}})$ is defined in \Cref{lem:PvimorphismUtoW}. 
    % This is a Galois extension and we have
    % $$\operatorname{Gal}\left(\C(\mathscr{U}_{\rm{VI}})/\C(\mathscr{W}_{\rm{VI}})\right)\cong W(D_4).$$
\end{proposition}

\begin{proof}
    We first establish the $\C(\mathscr{W}_{\rm{VI}})$-linear isomorphism of fields between $K_{\Gamma}$ and $\C(\mathscr{U}_{\rm{VI}})$. 
    To do this, we describe the incidence variety using (dual) Pl\"ucker coordinates on $\mathbb{G}(1,3)$ in the following way.
For a projective line in $\p^3$ defined as the intersection of two hyperplanes
\begin{equation*}
    \ell = V( a_0 X_0+a_1X_1+a_2X_2+a_3X_3, b_0 X_0+b_1X_1+b_2X_2+b_3X_3),
\end{equation*}
the Pl\"ucker coordinates of $\ell$ are $[P_{0,1}:P_{0,2}:P_{0,3}:P_{1,2}:P_{2,3}:P_{3:1}]\in V(P_{0,1}P_{2,3}+ P_{0,3}P_{1,2}+P_{0,2}P_{1,3})\subset \p^5$, defined by 
\begin{equation*}
    P_{i,j}=a_i b_j - a_j b_i, \quad i,j=0,\dots,3.
\end{equation*}

We work with the incidence variety as an algebraic set in $\mathscr{W}_{\rm{VI}}\times \p^5$ via Pl\"ucker coordinates of lines.
These can be computed from the expressions for lines in \eqref{eq:linesPVI} and are given in \Cref{tab:PVI:plucker:1to24}.
% The function field of the incidence variety is generated by the Pl\"ucker coordinates of the lines on $\overline{\mathcal{M}}_{w}$. 
In particular, by the last three columns of the table,
\begin{equation*}
    P_{2,3}P_{3,1}P_{1,2}=0
\end{equation*}
on $\Gamma$. Corresponding to the three components of $\Gamma$ defined by $P_{2,3}=0$, $P_{3,1}=0$ and $P_{1,2}=0$,
are respectively the sets of lines
 $\{L_1,\ldots, L_8\}$, $\{L_9,\ldots, L_{16}\}$ and $\{L_{17},\ldots, L_{24}\}$.
%  have Pl\"ucker coordinates contained respectively in the three components of $\Gamma$ defined by 
% $$P_{2,3}=0, \quad P_{3,1}=0, \quad P_{1,2}=0,$$
% each of which
Each component is a variety in $\mathscr{W}_{\rm{VI}}\times \p^5$ of dimension 4.

% \begin{equation}
%     \begin{array}{rlllllll}
%         L_1 :& 
%         &P_{0,1} = \frac{u_1}{u_2}+u_3+\frac{u_3 u_4}{u_1 u_2}+\frac{1}{u_4}, 
%         &P_{0,2} =-\frac{u_3^2}{u_2^2}-1, 
%         &P_{0,3} = -\frac{u_2}{u_3}-\frac{u_3}{u_2}, 
%         &P_{1,2} = \frac{u_3}{u_2},
%         &P_{2,3} = 0,  
%         &P_{3,1} = -1, \\
%         L_2 :& 
%         &P_{0,1} = \frac{u_2}{u_1}+\frac{u_1 u_2}{u_3 u_4}+u_4+\frac{1}{u_3}, 
%         &P_{0,2} =-\frac{u_2^2}{u_3^2}-1, 
%         &P_{0,3} = -\frac{u_2}{u_3}-\frac{u_3}{u_2}, 
%         &P_{1,2} = \frac{u_2}{u_3},
%         &P_{2,3} = 0,  
%         &P_{3,1} = -1, \\
%         L_3 :& 
%         &P_{0,1} = \frac{u_2}{u_1}+u_3+\frac{u_3 u_4}{u_1 u_2}+u_4, 
%         &P_{0,2} = -\frac{u_3^2 u_4^2}{u_1^2}-1, 
%         &P_{0,3} = -\frac{u_1}{u_3 u_4}-\frac{u_3 u_4}{u_1}, 
%         &P_{1,2} = \frac{u_3 u_4}{u_1},
%         &P_{2,3} = 0,  
%         &P_{3,1} = -1, \\
%         L_4 :& 
%         &P_{0,1} = \frac{u_1}{u_2}+\frac{u_2 u_1}{u_3 u_4}+\frac{1}{u_3}+\frac{1}{u_4}, 
%         &P_{0,2} = -\frac{u_1^2}{u_3^2 u_4^2}-1, 
%         &P_{0,3} = -\frac{u_1}{u_3 u_4}-\frac{u_3 u_4}{u_1}, 
%         &P_{1,2} = \frac{u_1}{u_3 u_4},
%         &P_{2,3} = 0,  
%         &P_{3,1} = -1,\\
%         L_5 :& 
%         &P_{0,1} = \frac{u_1}{u_2}+\frac{u_2 u_1}{u_3 u_4}+u_3+u_4, 
%         &P_{0,2} = -u_1^2-1, 
%         &P_{0,3} = -u_1-\frac{1}{u_1}, 
%         &P_{1,2} = u_1,
%         &P_{2,3} = 0,  
%         &P_{3,1} = -1,\\
%         L_6 :& 
%         &P_{0,1} = \frac{u_2}{u_1}+\frac{u_3 u_4}{u_1 u_2}+\frac{1}{u_3}+\frac{1}{u_4}, 
%         &P_{0,2} = -\frac{1}{u_1^2}-1, 
%         &P_{0,3} = -u_1-\frac{1}{u_1}, 
%         &P_{1,2} = \frac{1}{u_1},
%         &P_{2,3} = 0,  
%         &P_{3,1} = -1, \\
%         L_7:& 
%         &P_{0,1} = \frac{u_2}{u_1}+\frac{u_1 u_2}{u_3 u_4}+u_3+\frac{1}{u_4}, 
%         &P_{0,2} = -\frac{u_2^2}{u_4^2}-1, 
%         &P_{0,3} = -\frac{u_2}{u_4}-\frac{u_4}{u_2}, 
%         &P_{1,2} = \frac{u_2}{u_4},
%         &P_{2,3} = 0,  
%         &P_{3,1} = -1, \\
%         L_8:& 
%         &P_{0,1} = \frac{u_1}{u_2}+\frac{u_3 u_4}{u_1 u_2}+u_4+\frac{1}{u_3}, 
%         &P_{0,2} = -\frac{u_4^2}{u_2^2}-1, 
%         &P_{0,3} = -\frac{u_2}{u_4}-\frac{u_4}{u_2}, 
%         &P_{1,2} = \frac{u_4}{u_2},
%         &P_{2,3} = 0,  
%         &P_{3,1} = -1.
%     \end{array}
% \end{equation}



The lines $L_1,\dots,L_8$ correspond to points in the hypersurface $P_{2,3}=0$ with $P_{3,1}\neq0$, so to study this part of $\Gamma$ we restrict to the affine chart with coordinates $(p_{0,1},p_{0,2},p_{0,3},p_{1,2})\in\C^4$, 
$$[P_{0,1}:P_{0,2}:P_{0,3}:P_{1,2}:P_{2,3}:P_{3:1}] = [p_{0,1}:p_{0,2}:p_{0,3}:p_{1,2}:0:1].$$
The condition that $\ell\subset \overline{\mathcal{M}}_w$ and also the Pl\"ucker relation is written in these coordinates as
\begin{equation} \label{eq:pluckereqs1}
    \begin{gathered}
        p_{0,3} p_{0,2}^2+p_{1,2} p_{0,2}^2+p_{0,3}^2 p_{1,2} =0, \\
        w_3 p_{0,3} p_{1,2}^2 p_{0,2}-w_2 p_{0,3}^2 p_{1,2}^2-p_{0,1} p_{0,3} p_{0,2}^2-2 p_{0,1} p_{1,2} p_{0,2}^2=0,\\
        w_1 p_{0,3}^2 p_{1,2} p_{0,2}-w_3 p_{0,1} p_{0,3} p_{1,2} p_{0,2}+w_4 p_{0,3}^2 p_{1,2}^2+p_{0,1}^2 p_{0,2}^2+p_{0,3}^2 p_{0,2}^2=0,\\
        p_{0,2}+p_{0,3} p_{1,2} = 0.
    \end{gathered}
\end{equation}
All lines $L_1\dots,L_8$ have $p_{0,1},p_{0,2},p_{0,3},p_{1,2}\neq 0$, for generic $w$, and under this assumption the equations \eqref{eq:pluckereqs1} reduce, letting $p_{1,2}=-b$, to a single equation $\Lambda_1(b)=0$, where $\Lambda_1(b)$ is the degree eight polynomial
\begin{equation} \label{pvi:Lambda1}
    \Lambda_1(b)= b^8-w_1b^7 +w_4 b^6 + (w_1- w_2 w_3) b^5+ (w_2^2 +w_3^2-2w_4-2) b^4 +(w_1- w_2 w_3)b^3+ w_4 b^2- w_1 b+1.
\end{equation}
The Pl\"ucker coordinates of lines $L_9,\dots,L_{16}$ are given by the solutions of equations that similarly reduce to $\Lambda_2(b)=0$, where 
\begin{equation} \label{pvi:Lambda2}
    \Lambda_2(b)= b^8-w_2b^7 +w_4b^6 + (w_2 - w_1 w_3) b^5+(w_1^2+ w_3^2-2w_4-2)b^4+(w_2- w_1 w_3)b^3+ w_4 b^2 - w_2 b+1,
\end{equation}
while $L_{17},\dots,L_{24}$ lead to $\Lambda_3(b)=0$, where 
\begin{equation} \label{pvi:Lambda3}
    \Lambda_3(b)= b^8-w_3b^7 + w_4 b^6+ (w_3  - w_1 w_2) b^5+( w_1^2+ w_2^2-2  w_4-2)b^4 +(  w_3-w_1 w_2)b^3+ w_4 b^2 - w_3b+1.
\end{equation}

Recalling the definition of $b_k$, $1\leq k\leq 24$, in equation \eqref{eq:defbpvi}, the three polynomials \eqref{pvi:Lambda1}, \eqref{pvi:Lambda2}, \eqref{pvi:Lambda3} split over $\C(\mathscr{U}_{\rm{VI}})$, with roots $b_1,\dots,b_8$ for $\Lambda_1$,  $b_9,\dots,b_{16}$ for $\Lambda_2$, and $b_{17},\dots,b_{24}$ for $\Lambda_3$.
Therefore the irreducible components $\Gamma_1,\Gamma_2$, and $\Gamma_3$, corresponding to $L_1,\dots,L_8$, $L_9,\dots,L_{16}$, and  $L_{17},\dots,L_{24}$ respectively, have function fields whose normal closures are $\C(b_1,\dots,b_8)$, $\C(b_7,\dots,b_{16})$, and 
$\C(b_{17},\dots,b_{24})$ respectively.
This means that $K_{\Gamma}\cong\mathbb{C}(b_1,\dots, b_{24}) \cong \C(\mathscr{U}_{\rm{VI}})$.
%     An explicit isomorphism is given by
%     \begin{equation*}
%        \C(\mathscr{U}_{\rm{VI}})\rightarrow  \C(\Gamma),
%     \end{equation*}
% with
% \begin{equation*}
%     u_1\mapsto P_{3,1}P_{1,2} R_1+P_{2,3}P_{1,2}R_2+P_{2,3}P_{3,1}R_3
% \end{equation*}
%     % In terms of $\{u_1,u_2,u_3,u_4\}$, the Pl\"ucker coordinates of the lines $L_1,\dots,L_{24}$ are given in \Cref{tab:PVI:plucker:1to24}. It follows immediately from the data in the table that $\C(\Gamma)=\C(\mathscr{U}_{\rm{VI}})$.
    Further, since $\C(\mathscr{U}_{\rm{VI}
    })$ is the splitting field of the polynomial $\Lambda_1(b)\Lambda_2(b)\Lambda_3(b)$ over $\mathbb{C}(\mathscr{W}_{\rm{VI}})$, it is a Galois extension of $\mathbb{C}(\mathscr{W}_{\rm{VI}})$.

\begingroup

\setlength{\tabcolsep}{15pt} % Default value: 6pt
\renewcommand{\arraystretch}{1.75} % Default value: 1

\begin{table}[h]
    \begin{equation*}
    \begin{array}{c||c|c|c|c|c|c|}
             & P_{0,1}   & P_{0,2}    & P_{0,3}   & P_{1,2}   &  P_{2,3}   & P_{3,1}        \\
        \hline\hline 
        %%
        %%
        L_1  & \frac{u_1}{u_2}+u_3+\frac{u_3 u_4}{u_1 u_2}+\frac{1}{u_4}     & -\frac{u_3^2}{u_2^2}-1     & -\frac{u_2}{u_3}-\frac{u_3}{u_2}   & \frac{u_3}{u_2}  & 0   & -1           \\
        \hline
        %%
        L_2  &  \frac{u_2}{u_1}+\frac{u_1 u_2}{u_3 u_4}+u_4+\frac{1}{u_3}     & -\frac{u_2^2}{u_3^2}-1     & -\frac{u_2}{u_3}-\frac{u_3}{u_2}   & \frac{u_2}{u_3}  & 0   & -1           \\
        \hline
        %%
        L_3  & \frac{u_2}{u_1}+u_3+\frac{u_3 u_4}{u_1 u_2}+u_4     & -\frac{u_3^2 u_4^2}{u_1^2}-1  & -\frac{u_1}{u_3 u_4}-\frac{u_3 u_4}{u_1}  & \frac{u_3 u_4}{u_1}  & 0   & -1           \\
        \hline
        %%
        L_4  & \frac{u_1}{u_2}+\frac{u_2 u_1}{u_3 u_4}+\frac{1}{u_3}+\frac{1}{u_4}    &-\frac{u_1^2}{u_3^2 u_4^2}-1  & -\frac{u_1}{u_3 u_4}-\frac{u_3 u_4}{u_1}  & \frac{u_1}{u_3 u_4}  & 0   & -1           \\
        \hline
        %%
        L_5  & \frac{u_1}{u_2}+\frac{u_2 u_1}{u_3 u_4}+u_3+u_4    &-u_1^2-1  & -u_1-\frac{1}{u_1}  & u_1  & 0   & -1           \\
        \hline
        %%
        L_6  & \frac{u_2}{u_1}+\frac{u_3 u_4}{u_1 u_2}+\frac{1}{u_3}+\frac{1}{u_4}    &-\frac{1}{u_1^2}-1  & -u_1-\frac{1}{u_1}  & \frac{1}{u_1}  & 0   & -1           \\
        \hline
        %%
        L_7  & \frac{u_2}{u_1}+\frac{u_1 u_2}{u_3 u_4}+u_3+\frac{1}{u_4}    &-\frac{u_2^2}{u_4^2}-1  & -\frac{u_2}{u_4}-\frac{u_4}{u_2}  & \frac{u_2}{u_4}  & 0   & -1           \\
        \hline
        %%
        L_8  & \frac{u_1}{u_2}+\frac{u_3 u_4}{u_1 u_2}+u_4+\frac{1}{u_3}    &-\frac{u_4^2}{u_2^2}-1  & -\frac{u_2}{u_4}-\frac{u_4}{u_2}  & \frac{u_4}{u_2}  & 0   & -1           \\
        \hline\hline 
        L_9  & -\frac{u_3^2}{u_1^2}-1    &\frac{u_2}{u_1}+u_3+\frac{u_3 u_4}{u_1 u_2}+\frac{1}{u_4}  & -\frac{u_1}{u_3}-\frac{u_3}{u_1} & -\frac{u_3}{u_1}  & 1   & 0           \\
        \hline
        L_{10}  & -\frac{u_1^2}{u_3^2}-1 & \frac{u_1}{u_2}+\frac{u_2 u_1}{u_3 u_4}+u_4+\frac{1}{u_3} & -\frac{u_1}{u_3}-\frac{u_3}{u_1}  & -\frac{u_1}{u_3}  & 1   & 0           \\
        \hline
        L_{11} &-\frac{u_2^2}{u_3^2 u_4^2}-1    &\frac{u_2}{u_1}+\frac{u_1 u_2}{u_3 u_4}+\frac{1}{u_3}+\frac{1}{u_4}  & -\frac{u_2}{u_3 u_4}-\frac{u_3 u_4}{u_2} &-\frac{u_2}{u_3 u_4}  & 1   & 0           \\
        \hline
        L_{12}  & -\frac{u_3^2 u_4^2}{u_2^2}-1  &\frac{u_1}{u_2}+u_3+\frac{u_3 u_4}{u_1 u_2}+u_4  & -\frac{u_2}{u_3 u_4}-\frac{u_3 u_4}{u_2} & -\frac{u_3 u_4}{u_2} & 1   & 0          \\
        \hline
        L_{13}  & -\frac{1}{u_2^2}-1  &\frac{u_1}{u_2}+\frac{u_3 u_4}{u_1 u_2}+\frac{1}{u_3}+\frac{1}{u_4}  & -u_2-\frac{1}{u_2} & -\frac{1}{u_2}  & 1   & 0           \\
        \hline
        L_{14}  & -u_2^2-1  & \frac{u_2}{u_1}+\frac{u_1 u_2}{u_3 u_4}+u_3+u_4 & -u_2-\frac{1}{u_2} & -u_2  & 1   & 0           \\
        \hline
        L_{15}  & -\frac{u_4^2}{u_1^2}-1  & \frac{u_2}{u_1}+\frac{u_3 u_4}{u_1 u_2}+u_4+\frac{1}{u_3} & -\frac{u_1}{u_4}-\frac{u_4}{u_1} & -\frac{u_4}{u_1}  & 1   & 0           \\
        \hline
        L_{16}  & -\frac{u_1^2}{u_4^2}-1  & \frac{u_1}{u_2}+\frac{u_2 u_1}{u_3 u_4}+u_3+\frac{1}{u_4} & -\frac{u_1}{u_4}-\frac{u_4}{u_1}  & -\frac{u_1}{u_4} & 1   & 0          \\
        \hline \hline
        L_{17}  & -\frac{u_2^2}{u_1^2}-1  & -\frac{u_1}{u_2}-\frac{u_2}{u_1} &\frac{u_2}{u_3 u_4}+u_2+\frac{u_3}{u_1}+\frac{u_4}{u_1} & 0 & -1   & \frac{u_2}{u_1}         \\
        \hline
        L_{18}  &-\frac{u_1^2}{u_2^2}-1  & -\frac{u_1}{u_2}-\frac{u_2}{u_1} & \frac{u_1}{u_3}+\frac{u_1}{u_4}+\frac{u_3 u_4}{u_2}+\frac{1}{u_2}  & 0 & -1   & \frac{u_1}{u_2}         \\
        \hline
        L_{19}  & -u_4^2-1 & -u_4-\frac{1}{u_4} & \frac{u_1}{u_3}+u_2+\frac{u_3 u_4}{u_2}+\frac{u_4}{u_1} & 0 & -1   & u_4          \\
        \hline
        L_{20}  & -\frac{1}{u_4^2}-1  & -u_4-\frac{1}{u_4} & \frac{u_1}{u_4}+\frac{1}{u_2}+\frac{u_2}{u_3 u_4}+\frac{u_3}{u_1} & 0 & -1   & \frac{1}{u_4}          \\
        \hline
        L_{21}  & -u_3^2-1 & -u_3-\frac{1}{u_3} & \frac{u_1}{u_4}+u_2+\frac{u_3 u_4}{u_2}+\frac{u_3}{u_1} & 0 & -1   & u_3        \\
        \hline
        L_{22}  & -\frac{1}{u_3^2}-1   & -u_3-\frac{1}{u_3} & \frac{u_1}{u_3}+\frac{1}{u_2}+\frac{u_2}{u_3 u_4}+\frac{u_4}{u_1} & 0 & -1   & \frac{1}{u_3}          \\
        \hline
        L_{23}  & -\frac{u_1^2 u_2^2}{u_3^2 u_4^2}-1 & -\frac{u_1 u_2}{u_3 u_4}-\frac{u_3 u_4}{u_1 u_2} & \frac{u_1}{u_3}+\frac{u_1}{u_4}+u_2+\frac{u_2}{u_3 u_4}  & 0 & -1   & \frac{u_1 u_2}{u_3 u_4}          \\
        \hline
        L_{24}  & -\frac{u_3^2 u_4^2}{u_1^2 u_2^2}-1  & -\frac{u_1 u_2}{u_3 u_4}-\frac{u_3 u_4}{u_1 u_2} & \frac{u_4 u_3}{u_2}+\frac{u_3}{u_1}+\frac{u_4}{u_1}+\frac{1}{u_2} & 0 & -1   & \frac{u_3 u_4}{u_1 u_2}          \\
        \hline
    \end{array}
    \end{equation*}
    \caption{Pl\"ucker coordinates of lines $L_1,\dots,L_{24}$ on the monodromy surface for $\pain{VI}$}
    \label{tab:PVI:plucker:1to24}
\end{table}
\endgroup



Elements of $\operatorname{Gal}\left(\C(\mathscr{U}_{\rm{VI}})/\C(\mathscr{W}_{\rm{VI}})\right)$ give permutations of $b_1,\dots,b_{24}$, which in turn induce permutations of the lines $L_1,\dots,L_{24}$ which necessarily preserve incidence relations among them, yielding a homomorphism 
\begin{equation} \label{eq:galtofix}
\operatorname{Gal}\left(\C(\mathscr{U}_{\rm{VI}})/\C(\mathscr{W}_{\rm{VI}})\right)\rightarrow \operatorname{Fix}_{\operatorname{Aut}(\mathcal{G}_{\rm{VI}})}(\mathcal{G}_{\rm{VI}}^\infty).
\end{equation}
The kernel is trivial, since any nontrivial element of $\operatorname{Gal}\left(\C(\mathscr{U}_{\rm{VI}})/\C(\mathscr{W}_{\rm{VI}})\right)$ acts nontrivially on $\mathbb{C}(\mathscr{U}_{\rm{VI}})=\C(b_1,\dots,b_{24})$. 
On the other hand, the actions of $r_1,r_2,r_3,r_4$ on $\C(\mathscr{U}_{\rm{VI}})$ in \Cref{pvi:actionofD4onU} give elements of $\operatorname{Gal}\left(\C(\mathscr{U}_{\rm{VI}})/\C(\mathscr{W}_{\rm{VI}})\right)$
% \begin{equation}
% \begin{aligned}
%     g_1 &: &&u_1\to u_1, \quad 
%     &&u_2\to u_2, \quad 
%     &&u_3\to u_3, \quad 
%     &&u_4 \to \frac{u_1 u_2}{u_3 u_4}, \\
%     g_2 &: &&u_1\to \frac{u_4}{u_2}, \quad 
%     &&u_2\to \frac{u_4}{u_1}, \quad 
%     &&u_3\to \frac{u_3 u_4}{u_1 u_2}, \quad 
%     &&u_4\to u_4. \\
%     g_3 &: &&u_1\to \frac{u_3}{u_2}, \quad 
%     &&u_2\to \frac{u_3}{u_1}, \quad 
%     &&u_3 \to u_3, \quad 
%     &&u_4\to \frac{u_3 u_4}{u_1 u_2}, \\
%     g_4 &: &&u_1 \to u_1, \quad 
%     &&u_2\to \frac{u_1}{u_3}, \quad 
%     &&u_3\to \frac{u_1}{u_2}, \quad
%     &&u_4 \to u_4.
% \end{aligned}
% \end{equation}
which induce the permutations of lines as claimed, and we have the required isomorphism \eqref{eq:galtofix}.
\end{proof}

\end{document}
